
\documentclass[12pt,a4paper]{article}
\usepackage[T2A]{fontenc}
\usepackage[utf8]{inputenc}
\usepackage[russian]{babel}
\usepackage{float}
\usepackage{amsmath, amssymb}
\usepackage{geometry}
\usepackage{tikz}
\usepackage{tikz-3dplot}
\usepackage{caption}
\usetikzlibrary{arrows.meta, calc, decorations.pathmorphing}
\geometry{margin=2cm}

\newcommand{\va}{\bar a}
\newcommand{\vb}{\bar b}
\newcommand{\vc}{\bar c}
\newcommand{\vo}{\bar o}
\newcommand{\vx}{\bar x}
\newcommand{\vy}{\bar y}

\begin{document}

\setcounter{page}{101}

\subsection*{\textbf{ЛЕКЦИЯ 14. ЛИНЕЙНЫЕ ПРОСТРАНСТВА}}

\subsection*{\textbf{14.1. Понятие линейного пространства. Примеры
линейных пространств}}

\textbf{Определение.} Линейным пространством называется множество
$$
V=\{\va,\vb,\vc,\ldots\},
$$
где $\va,\vb,\vc,\ldots$ --- элементы пространства (векторы), если
для любых двух векторов $\va$ и $\vb$ и любого числа
$\lambda\in\mathbb R\Rightarrow \va+\vb\in V$ и $\lambda\va\in V$, и
выполняются следующие аксиомы:

1) $\va+\vb=\vb+\va$;

2) $(\va+\vb)+\vc=\va+(\vb+\vc)$ $\forall\va,\vb,\vc\in V$;

3) $\forall\va\in V$ $\exists$ нулевой вектор;
$$
\vo\in V:\quad \va+\vo=\vo+\va=\va;
$$

4) $\forall\va\in V$ $\exists$ противоположный вектор
$$
-\va\in V:\quad \va+(-\va)=(-\va)+\va=\vo;
$$

5) $(\lambda\cdot\mu)\va=\lambda(\mu\va)$;

6) $\lambda(\va+\vb)=\lambda\va+\lambda\vb$;

7) $(\lambda+\mu)\va=\lambda\va+\mu\va$ $\forall\va,\vb\in V$ и
$\forall\lambda,\mu\in\mathbb R$;

8) $1\cdot\va=\va$.

1) --- 4) --- аксиомы по операции сложения элементов, 5) --- 8) ---
аксиомы по операции умножения элемента на число.

\textbf{Замечание.} Линейное пространство называют также векторным
пространством.

\textbf{Примеры линейных пространств}

1) пространства геометрических векторов: $V_1$ --- на прямой; $V_2$
--- на плоскости; $V_3$ --- в пространстве;

2) множества $\mathbb R$ и $\mathbb C$;

3) множество матриц $\mathbb R^{m\times n}$;

4) множество многочленов степени $\le n$ --- $P_n$;

5) множество функций, непрерывных на отрезке $[a,b]$ --- $C[a,b]$;

6) $n$-мерное пространство арифметических векторов --- $\mathbb R^n$
--- множество всевозможных упорядоченных наборов из $n$
действительных чисел, называемых арифметическими векторами.
$$
\mathbb R^n=\{\vx=(x_1,x_2,\ldots,x_n);\ x_i\in\mathbb R\}
$$
Пусть $\vx=(x_1,x_2,\ldots,x_n)$ и $\vy=(y_1,y_2,\ldots,y_n)$.

Тогда:

1) $\vx=\vy\Leftrightarrow x_i=y_i$, $\forall i=1,2,\ldots,n$;

2) $\vx+\vy=(x_1+y_1,\ldots,x_n+y_n)$;

3) $\lambda\vx=(\lambda x_1,\lambda x_2,\ldots,\lambda x_n)$.

\subsection*{\textbf{14.2. Простейшие свойства линейных пространств}}

1) В линейном пространстве существует единственный нулевой элемент $\vo$.

\textbf{Доказательство}

Пусть $\exists\bar o_1$ и $\bar o_2$ --- два различных нулевых
вектора. Тогда из аксиомы 3) следует, что $\bar o_1+\bar o_2=\bar o_1=\bar o_2$.

\textbf{ч.т.д.}

2) Для любого вектора линейного пространства существует единственный
противоположный вектор.

\textbf{Доказательство}

Пусть $\va'$ и $\va''$ --- два различных противоположных элемента к
элементу $\va$.

Тогда $\va+\va'=\va'+\va=\vo$ и $\va+\va''=\va''+\va=\vo$.
$$
\Rightarrow \va'=\va'+\vo=\va'+(\va+\va'')=(\va'+\va)+\va''=
$$
$$
=\vo+\va''=\va''.
$$
\textbf{ч.т.д.}

3)
\begin{equation}
0\cdot\va=\vo \tag{14.1}
\end{equation}
и
\begin{equation}
\lambda\cdot\vo=\vo. \tag{14.2}
\end{equation}

\textbf{Доказательство}

Для доказательства равенства (14.1) достаточно доказать, что
$\vb+0\cdot\va=\vb$, $\forall\vb\in V$.
$$
\vb+0\cdot\va=(\vb+\vo)+0\cdot\va=\vb+((-\va)+\va)+0\cdot\va=
$$
$$
=(\vb+(-\va))+1\va+0\va=(\vb+(-\va))+(1+0)\cdot\va=
$$
$$
=(\vb+(-\va))+\va=\vb+((-\va)+\va)=\vb+\vo=\vb.
$$
Равенство (14.2) доказывается, используя равенство (14.1) и аксиому 5):
$$
\lambda\cdot\vo=\lambda(0\cdot\va)=(\lambda\cdot0)\va=0\cdot\va=\vo
$$
\textbf{ч.т.д.}

4) В линейном пространстве из равенства $\lambda\va=\vo$ следует, что
либо $\lambda=0$, либо $\va=\vo$.

\textbf{Доказательство}

Из равенства (14.1) следует, что случай $\lambda=0$ возможен, если
$\lambda\cdot\va=\vo$.

Если $\lambda\ne0$, то
$$
\va=1\cdot\va=\left(\frac1\lambda\lambda\right)\va=\frac1\lambda(\lambda\va)=\frac1\lambda\cdot\vo=\vo.
$$
\textbf{ч.т.д.}

5) $-\va=(-1)\va$

\textbf{Доказательство}
$$
\va+(-1)\va=1\va+(-1)\va=(1-1)\va=0\cdot\va=\vo.
$$
\textbf{ч.т.д.}

6) Для любых двух векторов $\va$ и $\vb$ в линейном пространстве
существует, и притом единственная, разность $\vb-\va$.

\textbf{Доказательство}
$$
\va+(\vb+(-\va))=(\va+(-\va))+\vb=\vo+\vb=\vb\Rightarrow \vb-\va=\vb+(-\va);
$$
Докажем единственность. Пусть $\exists\vc=\vb-\va$ --- другая разность, но
$$
\vc=\vc+\vo=\vc+(\va+(-\va))=
$$
$$
=(\vc+\va)+(-\va)=\vb+(-\va).
$$
\textbf{ч.т.д.}

\subsection*{\textbf{14.3. Линейное подпространство}}

\textbf{Определение.} Подмножество $M\subset V$ называется линейным
подпространством линейного пространства $V$, если оно само является
линейным пространством относительно операции сложения элементов и
умножения элемента на число.

\textbf{Замечание.} Подмножество $M\subset V$ есть линейное
подпространство в линейном пространстве $V$ тогда и только тогда, когда

1) $\forall\va,\vb\in M\Rightarrow \va+\vb\in M$

2) $\forall\lambda\in\mathbb R$ и $\forall\va\in M\Rightarrow \lambda\va\in M$

\textbf{Замечание.} Подмножество $M\subset V$ есть линейное
подпространство в линейном пространстве $V$ тогда и только тогда, когда
$$
\forall\va,\vb\in M \text{ и } \forall\alpha,\beta\in\mathbb
R\Rightarrow \alpha\va+\beta\vb\in M
$$

\textbf{Пример}

Доказать, что множество геометрических векторов $V^3$ --- линейное пространство.

\textbf{Доказательство}

$\forall\va,\vb\in V^3$ и $\forall\lambda\in\mathbb R$
$$
\Rightarrow \va+\vb\in V^3 \text{ и } \lambda\va\in V^3.
$$

1) $\va+\vb=\vb+\va$ $\forall\va,\vb\in V^3$;

2) $(\va+\vb)+\vc=\va+(\vb+\vc)$ $\forall\va,\vb,\vc\in V^3$;

3) Для $\forall\va\in V^3$ $\exists\vo\in V^3$: $\va+\vo=\vo+\va=\va$;

4) Для $\forall\va\in V^3$ $\exists-\va\in V^3$: $\va+(-\va)=(-\va)+\va=\vo$;

5) $(\lambda\cdot\mu)\va=\lambda(\mu\va)$ $\forall\va\in V^3$ и
$\forall\lambda,\mu\in\mathbb R$;

6) $\lambda(\va+\vb)=\lambda\va+\lambda\vb$ $\forall\va,\vb\in V^3$ и
$\forall\lambda\in\mathbb R$;

7) $(\lambda+\mu)\va=\lambda\va+\mu\va$ $\forall\va\in V^3$ и
$\forall\lambda,\mu\in\mathbb R$;

8) $1\cdot\va=\va$ $\forall\va\in V^3$.

$\Rightarrow V^3$ --- линейное пространство

\textbf{Пример}

Доказать, что множество матриц размера $2\times2$ --- $\mathbb
R^{2\times2}$ является линейным пространством.

\textbf{Доказательство}

$\forall A,B\in\mathbb R^{2\times2}$ и $\forall\lambda\in\mathbb R$
$$
\Rightarrow A+B\in\mathbb R^{2\times2} \text{ и } \lambda A\in\mathbb
R^{2\times2}.
$$

1) $A+B=B+A$ $\forall A,B\in\mathbb R^{2\times2}$;

2) $(A+B)+C=A+(B+C)$ $\forall A,B,C\in\mathbb R^{2\times2}$;

3) Для $\forall A\in\mathbb R^{2\times2}$ $\exists O\in\mathbb
R^{2\times2}$: $A+O=O+A=A$;

4) Для $\forall A\in\mathbb R^{2\times2}$ $\exists -A\in\mathbb
R^{2\times2}$: $A+(-A)=(-A)+A=O$;

5) $(\lambda\cdot\mu)A=\lambda(\mu A)$ $\forall A\in\mathbb
R^{2\times2}$ и $\forall\lambda,\mu\in\mathbb R$;

6) $\lambda(A+B)=\lambda A+\lambda B$ $\forall A,B\in\mathbb
R^{2\times2}$ и $\forall\lambda\in\mathbb R$;

7) $(\lambda+\mu)A=\lambda A+\mu A$ $\forall A\in\mathbb
R^{2\times2}$ и $\forall\lambda,\mu\in\mathbb R$;

8) $1\cdot A=A$ $\forall A\in\mathbb R^{2\times2}$.

$\Rightarrow\mathbb R^{2\times2}$ --- линейное пространство.

\textbf{Пример}

Пусть $M$ --- множество многочленов степени $\le4$, которые делятся
на $t^2+t+1$. Доказать, что $M$ --- линейное подпространство в
линейном пространстве $P_4$.

\textbf{Доказательство}
$$
M=\{p(t)=(at^2+bt+c)(t^2+t+1);\ a,b,c\in\mathbb R\}
$$
$$
p(t)=(at^2+bt+c)(t^2+t+1)=
$$
$$
=a(t^4+t^3+t^2)+b(t^3+t^2+t)+c(t^2+t+1).
$$
Пусть
$$
p_1(t)=a_1(t^4+t^3+t^2)+b_1(t^3+t^2+t)+c_1(t^2+t+1);
$$
$$
p_2(t)=a_2(t^4+t^3+t^2)+b_2(t^3+t^2+t)+c_2(t^2+t+1);
$$
$$
\alpha p_1(t)+\beta p_2(t)=\alpha a_1(t^4+t^3+t^2)+\alpha b_1(t^3+t^2+t)+
$$
$$
+\alpha c_1(t^2+t+1)+\beta a_2(t^4+t^3+t^2)+\beta b_2(t^3+t^2+t)+
$$
$$
+\beta c_2(t^2+t+1)=(\alpha a_1+\beta a_2)(t^4+t^3+t^2)+
$$
$$
+(\alpha b_1+\beta b_2)(t^3+t^2+t)+(\alpha c_1+\beta c_2)(t^2+t+1)=
$$
$$
=\left[
\begin{array}{l}
\alpha a_1+\beta a_2=a\\
\alpha b_1+\beta b_2=b\\
\alpha c_1+\beta c_2=c
\end{array}
\right]=a(t^4+t^3+t^2)+b(t^3+t^2+t)+c(t^2+t+1)
$$
$$
\Rightarrow \alpha p_1(t)+\beta p_2(t)\in M
$$
$\forall p_1(t),p_2(t)\in M$ и $\forall\alpha,\beta\in\mathbb R$.

$\Rightarrow M$ --- линейное подпространство в линейном пространстве $P_4$

\textbf{Пример}

Пусть $M$ --- множество симметричных матриц размера $2\times2$.
Доказать, что $M$ --- линейное подпространство в линейном
пространстве $\mathbb R^{2\times2}$.

\textbf{Доказательство}
$$
M=\left\{X=
\begin{pmatrix}a&c\\ c&b
\end{pmatrix};\ a,b,c\in\mathbb R\right\}.
$$
Пусть
$$
A=
\begin{pmatrix}a_1&c_1\\ c_1&b_1
\end{pmatrix};\qquad
B=
\begin{pmatrix}a_2&c_2\\ c_2&b_2
\end{pmatrix}.
$$
$$
\alpha A+\beta B=
\begin{pmatrix}\alpha a_1&\alpha c_1\\ \alpha c_1&\alpha b_1
\end{pmatrix}+
\begin{pmatrix}\beta a_2&\beta c_2\\ \beta c_2&\beta b_2
\end{pmatrix}=
$$
$$
=
\begin{pmatrix}
\alpha a_1+\beta a_2&\alpha c_1+\beta c_2\\
\alpha c_1+\beta c_2&\alpha b_1+\beta b_2
\end{pmatrix},
$$
$$
\Rightarrow \alpha A+\beta B\in M \quad \forall A,B\in M \text{ и }
\forall\alpha,\beta\in\mathbb R
$$
$$
\Rightarrow M \text{ --- линейное подпространство в линейном
пространстве } \mathbb R^{2\times2}.
$$

\textbf{Пример}

Пусть $M$ --- это множество функций вида
$$
f(t)=a+b\cdot\cos^2 t+c\cdot\cos2t.
$$
Доказать, что $M$ --- линейное пространство.

\textbf{Доказательство}
$$
f(t)=\alpha+\beta\cdot\cos^2 t+\gamma\cdot\cos2t=
\alpha+\beta\frac{1+\cos2t}{2}+\gamma\cos2t=
$$
$$
=\alpha+\frac{\beta}{2}+\frac{\beta}{2}\cos2t+\gamma\cos2t=\left(\alpha+\frac{\beta}{2}\right)+\left(\frac{\beta}{2}+\gamma\right)\cos2t=
$$
$$
=\left[
\begin{array}{l}
\alpha+\dfrac{\beta}{2}=a\\[1mm]
\dfrac{\beta}{2}+\gamma=b
\end{array}
\right]=a+b\cos2t,
$$
$$
\Rightarrow M=\{f(t)=a+b\cos2t;\ a,b\in\mathbb R\}
$$
Рассмотрим $C(-\infty;+\infty)$ --- линейное пространство функций,
непрерывных на интервале $(-\infty;+\infty)$.
$$
M\subset C(-\infty;+\infty)
$$
Докажем, что $M$ --- линейное подпространство в линейном пространстве
$C(-\infty;+\infty)$.

Пусть
$$
f_1(t)=a_1+b_1\cos2t;
$$
$$
f_2(t)=a_2+b_2\cos2t.
$$
$$
\lambda f_1(t)+\mu f_2(t)=\lambda a_1+\lambda b_1\cos2t+\mu a_2+\mu b_2\cos2t=
$$
$$
=(\lambda a_1+\mu a_2)+(\lambda b_1+\mu b_2)\cos2t=
\left[
\begin{array}{l}
\lambda a_1+\mu a_2=a\\
\lambda b_1+\mu b_2=b
\end{array}
\right]=a+b\cos2t
$$
$$
\Rightarrow \lambda f_1(t)+\mu f_2(t)\in M
$$
$\forall f_1(t),f_2(t)\in M$ и $\forall\lambda,\mu\in\mathbb R$
$$
\Rightarrow M \text{ --- линейное подпространство в линейном
пространстве } C(-\infty;+\infty).
$$
Т.е. $M$ --- само по себе является линейным пространством.

\end{document}
